\section[Naissance de la bibliothèque]{Une bibliothèque née de l'université de droit}
L'histoire de la BIU Cujas remonte au XVe siècle où l'on mentionne pour la première fois une  \enquote{bibliothèque de l'Ecole de droit} en 1475. Cujas continua d'évoluer jusqu'à s'imposer devant une bibliothèque de référence en droit grâce à \enquote{la richesse de ses collections, l'étendue des services proposés, les missions particulières qui lui sont attribuées, ainsi que l'expérience acquise par son personnel}.

En 1806, la Faculté de droit de Paris réouvre ses portes, mais c'est en 1829 qu'ouvre officiellement la bibliothèque Cujas. Elle est crée comme bibliothèque de la Faculté de droit de Paris, héritière de la bibliothèque de l'Ecole de droit de l'Université de paris. En 1829, elle permet aux étudiants de préparer leurs thèses et autorise en 1831 tous les étudiants de droit de la capitale à pouvoir consulter ces collections. En 1876, Paul Viollet est nommé comme le premier bibliothécaire professionnel à la tête de la bibliothèque. Sous sa direction, entre 1876 et 1914, la bibliothèque Cujas connaît une forte expansion et un développement de ses collections, une augmentation de son budget et l'accroissement de son effectif.   

En 1970, la bibliothèque Cujas est confronté comme les autres bibliothèques aux conséquences de la partition de l'Université de Paris, entraînant la disparition de l'ancienne faculté de droit. Depuis la signature du traité de paix, chacun pense que \enquote{l'on va pouvoir beaucoup changer l'état des choses}, mais en réalité les années 1919-1939 sont \enquote{une période noire pour les bibliothèques d'études et de recherche}. Le contexte n'est pas favorable au développement des bibliothèques, ni à une refonte de nos institutions publiques à cause de l'effondrement de la Troisième République. A la sortie de la Seconde Guerre Mondiale, la Société des Nations a conscience de l'importance que revêt la culture, crée dans une \enquote{Commission internationale de coopération intellectuelle, présidée par Henri Bergson}. Toutefois, l'appui financier venant à manquer et la multiplication des bibliothèques et des étudiants rendent difficile la mise en place des réformes, ainsi que son organisation. \enquote{Depuis le décret du 9 novembre 1946, il n'existe plus une bibliothèque de l'université de Paris mais des bibliothèques facultaires, auxquelles sont théoriquement rattachées les bibliothèques de laboratoires et dinstitutions de la faculté correspondante}. Après mai 1968 et la loi Faure, afin d'éviter la dispersion excessive des documents, il a été décidé de constituer des bibliothèques interuniversitaires. En effet, ne souhaitant pas voir disperser les fonds particuliers qui ont une importance particulière pour la recherche, \enquote{un rôle nationale, rendait difficile leur attribution à une université plutôt qu'à une autre}. C'est dans ce contexte que sont crées les bibliothèques interuniversitaires. La bibliothèque Cujas se trouve \enquote{intégrée dans la bibliothèque interuniversitaire A}. A ce moment-là, elle change de statut pour devenir la Bibliothèque universitaire Cujas. En 1979, la BIU Cujas est rattaché administrativement à l'Université Paris I, Panthéon-Sorbonne. Puis, grâce à une convention signée cette même année, elle se trouve co-gérée et conventionnée entre les universités Panthéon-Sorbonne (Paris I) et Panthéon-Assas(Paris II).

La BIU Cujas est spécialisé en droit, économmie et sciences politiques devenue aujourd'hui une \enquote{bibliothèque de référence nationale et internationale dans ces domaines}.
L'une des missions de Cujas est \enquote{d'offrir à la communauté universitaire et scientifique l'exhaustivité des ressources documentaires françaises ainsi qu'une vaste sélection de la production étrangère en sciences juridiques, économiques et politiques}. Le renom de Cujas vient de sa richesse documentaire \enquote{d'environ un million de documents (monographies, encyclopédies, ouvrages à feuillets mobiles, périodiques, thèses, mémoires) dont de nombreux ouvrages précieux}. Cujas conserve à l'intérieur de ses murs \enquote{sept incunables, le plus ancien datant de 1472}.

L'un des défis majeurs auxquels doit faire face la bibliothèque est la saturation de ces magasins, un problème auquel les institutions patrimoniales sont confrontées régulièrement. Les collections de Cujas augmentant chaque année, commencent à poser un problème de place dans ses magasins déjà saturés.